% Modernised reading — fonds Alexandre Grothendieck, shelfmark 117
% (pages 1-25, the whole folder).
%
% This is an INTERPRETATION, not a transcription. It re-reads the pages in
% today's notation and vocabulary, states the hypotheses the manuscript leaves
% implicit, and drops the critical apparatus so that the mathematics reads
% straight through. Everything the brackets and underlines carried moves into a
% footnote. It is held to being mathematically correct as it stands; where the
% manuscript is loose or mistaken, the true statement is what appears here and
% the footnote says what the page has. In any dispute of substance the
% transcription governs: whoever cites Grothendieck must cite batch-01.fr.tex
% (pp. 1-20) or batch-02.fr.tex (pp. 21-25).
%
% Pass: Opus 5.5 (claude-opus-5-5), 2026-09-23 — first pass, unchecked against the pages by a human.
% The folder was read whole, its two batches together; this is one document
% for the shelfmark. It was made on Opus 5.5 and has not been measured against
% the reference reading of folder 115 (made on Opus 5). The literature named
% in the footnotes is cited from memory and was not checked against the
% sources.
%
% Scope. Pages 6, 8, 10, 12, 14, 16, 18 (versos: administrative tables in
% another hand) and 23, 25 (typed university notices) carry nothing of his and
% are absent. Two leaves carry dates in his hand: « déc. 83 » (p. 1) and
% « 24.11.83 » (p. 4, margin); nothing else dates a leaf, and nothing orders
% the runs among themselves.
%
% Spine. What a « homotopy theory » attached to every small category should
% be — today a (pre)derivator — in four stretches:
%   pp. 1-2    « Structure des dérivateurs »: three levels of Hom (set, space,
%              internal space) and the adjunctions f_! -| f^* -| f_* at each.
%   pp. 3-4    a programme a)-i): models, long exact sequences, resolutions,
%              mapping spaces, invariants, K-theory, internal Hom; « tout …
%              en termes de (M, W) ».
%   pp. 5-20   his paginated run 1-9: the 2-functor A -> H(A), (Hot 1)
%              sums to products, a cancelled first integration through
%              presheaves (p. 7), (Hot 2)-(Hot 3) the adjoints f_!, f_*,
%              duality, the additive case; examples 1-4 (diagrams in M; W^{-1}
%              of diagrams in (M, W); Cat with a fundamental localizer;
%              derived categories of sheaves); (Hot 4)-(Hot 6).
%   pp. 21-24  cofinality when i has a left adjoint; a gluing condition for
%              A = A' u A''; the staircase of the loop (fibre) sequence.
%
% Notation fixed for the document. H for his script H (the theory), H(A) its
% value; f^*, f_!, f_* without his index H; Hot_H = H(Δ_0), Δ_0 the point
% category (his notation, kept); ∫_A = (π_A)_! (his « intégration », a
% homotopy colimit) and ∏_A = (π_A)_* (his « cointégration », a homotopy
% limit); A^ the presheaf category, h_A the Yoneda functor. « Cofibré » and
% « fibré » for Grothendieck's cofibred / fibred categories (op-fibration /
% fibration), read as what the page's « cofibrant », « fibrant » mean on p. 20.
%
% Substantive departures, each footnoted where it occurs:
%  p. 1    « Homhot(x, y) ≃ ∫_I Homhot(x, y) »: the global mapping space is a
%          homotopy LIMIT (his ∏ of p. 11), not ∫; the internal formulas
%          (3'), (3'') are stated as the page states them, the internal
%          object not being defined there.
%  p. 9    f^H_! : H(B) -> H(A) written as f^* (sic), corrected.
%  p. 13   (Hot 1) for Ex. 2: a localisation need not commute with an
%          infinite product; holds for model categories (cited from memory).
%  p. 17   the Remarque needs 2-out-of-3 for W' (supplied), then units and
%          counits are in W, W'.
%  p. 20   Ex. 4: the page denies a left adjoint in general, then derives
%          L f_! from an « exact à gauche » f_!; a left Kan extension of
%          abelian sheaves is right exact, not exact, so the argument fails
%          as written; both adjoints exist on the derived categories
%          (derivator of a Grothendieck abelian category, cited from memory).
%  p. 20   (Hot 4): the condition belongs on f' (f inherits it); the page's
%          bracketed alternatives « ou g fibrant ?? » are right (proper /
%          smooth functors, cited from memory).
%  p. 20   (Hot 6) is FALSE for a localisation of categories in the ordinary
%          sense: a category whose nerve is S^2, localised at all arrows,
%          gives Δ_0, and f^* is not fully faithful on spaces
%          (counterexample ours); true for localisations in the homotopical
%          sense.
%  p. 21   the page writes i where the argument needs j (p_0 j = p); j.
%  p. 24   the squares are not said to be homotopy cartesian; supplied.
%
% Announced and never established: the existence of the f_* for a
% fundamental localizer (« non encore »); the comparison of two (M, W)
% giving one theory (p. 19); the axioms of the programme (« Dégager les bons
% axiomes »); the gluing condition of p. 22 (wished for, « on veut »); the
% filtered case (p. 21 margin).
\documentclass[11pt,a4paper]{article}
% Le préambule commun à toutes les transcriptions du fonds Grothendieck.
%
% Il tient en une idée : **l'incertitude se compose, elle ne se résout pas.**
% Une transcription qui lisse un mot douteux a détruit la seule chose pour
% laquelle on la faisait — un lecteur doit pouvoir savoir, à chaque endroit,
% ce qui était sur le feuillet et ce qui a été ajouté. Les six macros qui
% suivent sont donc l'appareil critique, et non de la décoration.
%
% Les mêmes commandes sont reconnues par `scripts/render.mjs`, qui produit la
% vue de lecture du site. Une macro ajoutée ici doit l'être aussi là-bas,
% faute de quoi le rendu échouera bruyamment — ce qui est voulu : un rendu
% silencieusement faux est pire qu'un rendu absent.

\ProvidesPackage{grothendieck}[2026/08/08 Transcriptions du fonds Grothendieck]

\usepackage{amsmath,amssymb,amsthm}
\usepackage{mathrsfs}
\usepackage{stmaryrd}
% Les diagrammes commutatifs sont partout dans ce fonds : c'est la langue
% même de la géométrie algébrique de Grothendieck, et une transcription qui
% les rendrait en prose ne transcrirait rien.
\usepackage{tikz-cd}
% Français par défaut — c'est la langue des feuillets — mais l'anglais est
% chargé aussi : la traduction et le résumé sont des documents anglais, et la
% typographie française (espaces avant les deux-points) y serait une faute.
% Ces documents basculent par \selectlanguage{english} en tête de corps.
\usepackage[english,french]{babel}
\usepackage{fontspec}
\usepackage{geometry}
\geometry{a4paper,margin=25mm,marginparwidth=28mm}
\usepackage{marginnote}
\usepackage{xcolor}
% ulem, not soul. `soul`'s \st and \ul are built on a character-by-character
% scan that XeLaTeX's font machinery breaks: the document fails with an
% undefined control sequence rather than degrading. ulem does the same two jobs
% and is Unicode-engine safe. `normalem` stops it hijacking \emph, which the
% transcriptions use for Grothendieck's own emphasis.
\usepackage[normalem]{ulem}
\usepackage{hyperref}
\hypersetup{colorlinks=true,linkcolor=black,urlcolor=black,pdfborder={0 0 0}}

% --- Métadonnées du lot -----------------------------------------------------
% Reprises telles quelles par le script de rendu, qui les affiche en tête de
% la vue de lecture. Elles fixent ce que le fichier prétend couvrir : sans
% elles, une transcription flotte sans point d'attache dans un fonds de seize
% mille feuillets.

\newcommand{\folder}[1]{\gdef\grothFolder{#1}}
\newcommand{\batch}[1]{\gdef\grothBatch{#1}}
\newcommand{\leaves}[2]{\gdef\grothFirst{#1}\gdef\grothLast{#2}}
\newcommand{\foldertitle}[1]{\gdef\grothFoldertitle{#1}}
\gdef\grothFolder{?}\gdef\grothBatch{?}\gdef\grothFirst{?}\gdef\grothLast{?}\gdef\grothFoldertitle{}

% --- Le résumé --------------------------------------------------------------
%
% Il ouvre la lecture modernisée, sous le titre « Résumé », et il est composé
% en petit. Ce n'est pas un ornement : c'est le seul passage du document qui
% ne soit pas des mathématiques, et un lecteur doit pouvoir le reconnaître
% avant de l'avoir lu — puis le sauter s'il connaît déjà le sujet, ou s'y
% arrêter s'il ne le connaît pas. Le filet qui le suit marque la reprise du
% corps.

\newenvironment{resume}
  {\small\setlength{\parskip}{0.45em}}
  {\par\medskip\hrule\medskip}

% --- Mots-clés ---------------------------------------------------------------
%
% En anglais, en clôture du résumé de la lecture modernisée : le vocabulaire
% moderne sous lequel chercher ce que les feuillets construisent. Le manifeste
% du site (`npm run manifest`) les extrait pour étiqueter la cote entière —
% cette ligne est la source unique des tags, il n'y a pas de fichier à côté.
\newcommand{\keywords}[1]{\par\smallskip\noindent
  {\footnotesize\textsc{Keywords} --- \textit{#1}}\par}

% --- L'appareil critique ----------------------------------------------------

% Le numéro de feuillet, en marge. C'est l'ancre qui permet de régler un
% désaccord sur une formule en regardant *un* feuillet plutôt qu'une chemise
% de six cents. Le site s'en sert aussi pour tourner les pages du fac-similé
% quand on fait défiler la transcription.
\newcommand{\leaf}[1]{%
  \marginnote{\footnotesize\color{gray!70}\textsf{#1}}%
  \label{leaf:#1}\ignorespaces}

% Illisible : on ne devine pas. Un mot inventé qui se lit comme les autres est
% le pire résultat possible.
\newcommand{\ill}{\textcolor{red!60!black}{[\,\ldots\,]}}

% Lecture douteuse : le mot est donné, et signalé comme incertain.
\newcommand{\uncertain}[1]{\uline{#1}}

% Ajout de l'éditeur — une lettre manquante, un mot sous-entendu.
\newcommand{\add}[1]{\textcolor{blue!50!black}{[#1]}}

% Biffé par Grothendieck lui-même. Ce qu'il a rayé fait partie du document :
% une rature dit souvent où la pensée a changé de direction.
\newcommand{\struck}[1]{\sout{#1}}

% Note du transcripteur, sur l'état matériel du feuillet.
\newcommand{\note}[1]{\par\smallskip{\small\color{gray!60!black}\textit{[#1]}}\par\smallskip}

% Note marginale de Grothendieck, distincte du corps du texte.
\newcommand{\marginal}[1]{\par\smallskip\noindent\hspace{1em}%
  {\small\color{gray!60!black}#1}\par\smallskip}

% --- Titre ------------------------------------------------------------------

\AtBeginDocument{%
  \begin{center}
    {\large\bfseries Fonds Alexandre Grothendieck --- cote n\textsuperscript{o}~\grothFolder}\\[2pt]
    {\itshape\grothFoldertitle}\\[4pt]
    {\small feuillets \grothFirst--\grothLast\ (lot \grothBatch)}\\[2pt]
    {\footnotesize\color{gray!60!black}Transcription établie d'après le fac-similé numérisé
     par l'Université de Montpellier.\\
     Les lectures incertaines sont soulignées, les illisibles notées [\,\ldots\,],
     les ajouts entre crochets.}
  \end{center}
  \vspace{1em}}

\endinput


\folder{117}
\pages{1}{25}
\dating{1983}
\watermark{Édition de démonstration\\interprétation personnelle de l'œuvre}
\foldertitle{Théories homotopiques : notes manuscrites (1983, s.d.) — lecture modernisée du dossier entier}

\begin{document}

\section*{Résumé}

\begin{resume}
En topologie, on ne regarde souvent les espaces qu'« à déformation près » :
deux espaces qu'on peut déformer continûment l'un en l'autre comptent pour le
même. On obtient ainsi une catégorie, la catégorie homotopique, où beaucoup
de constructions familières se dérobent. Un recollement, une limite, un
quotient ne s'y calculent plus bien : deux diagrammes d'espaces peuvent être
indiscernables objet par objet et donner, une fois recollés, des résultats
différents. Le remède connu consiste à ne jamais perdre de vue les
diagrammes eux-mêmes. Il ne suffit pas de garder la catégorie homotopique des
espaces ; il faut garder, pour chaque forme de diagramme, la catégorie
homotopique des diagrammes de cette forme, et les opérations qui passent de
l'une à l'autre.

Ces feuillets cherchent à dire ce qu'est, en général, une telle « théorie
homotopique ». À chaque petite catégorie $A$ --- une forme de diagramme ---
elle associe une catégorie $\mathcal H(A)$ ; à chaque foncteur entre formes,
un moyen de restreindre, et deux moyens de pousser en avant, l'un qui recolle
(une « intégration », qui généralise la limite inductive), l'autre qui
rassemble (une « cointégration », qui généralise la limite projective). On
peut penser à une carte routière dessinée à plusieurs échelles : ce qui
compte n'est pas une carte, mais le système des cartes et les règles pour
passer d'une échelle à l'autre. Grothendieck pose des axiomes, les numérote
(Hot~1) à (Hot~6), les éprouve sur quatre exemples --- diagrammes dans une
catégorie ordinaire, diagrammes dans une catégorie munie d'une classe de
« flèches à inverser », les petites catégories elles-mêmes, les complexes de
faisceaux --- et note ce qu'il n'a pas encore su démontrer.

Deux feuillets sont datés de sa main, de novembre et de décembre 1983 ; le
feuillet 1, celui de décembre, porte le titre « Structure des
dérivateurs ». D'autres dressent un programme --- modèles, suites exactes, foncteurs dérivés, invariants --- qui
finit sur une consigne : « Dégager les bons axiomes ». Les derniers essaient
des énoncés plus fins : quand l'intégration sur une sous-catégorie donne la
même chose que sur le tout, comment une théorie se recolle sur deux morceaux,
et comment se déroule, en escalier, la suite des espaces de lacets.

Ce qu'on appelle aujourd'hui un \emph{dérivateur}, ou, sous d'autres
axiomes, une « théorie homotopique » au sens de Heller, est exactement ce
genre d'objet. Deux énoncés du dossier sont à corriger, et on le fait en
note : l'axiome (Hot~6) est faux pour une localisation au sens ordinaire, et
un argument sur les faisceaux ne tient pas tel qu'il est écrit.

\keywords{derivator, prederivator, homotopy Kan extension, homotopy exact square, homotopical category, fundamental localizer, fiber sequence}
\end{resume}

\section*{Le fil du dossier, et les conventions}

\begin{enumerate}
\item[1.] \emph{La structure des dérivateurs} (pages 1 et 2) : trois
niveaux de $\mathrm{Hom}$ et les adjonctions à chaque niveau.
\item[2.] \emph{Un programme} (pages 3 et 4), en neuf points.
\item[3.] \emph{Les axiomes} (pages 5 à 11), dans une suite qu'il pagine de 1
à 4 : la variance, (Hot~1) à (Hot~3), la dualité, le cas additif.
\item[4.] \emph{Quatre exemples} (pages 13 à 20, ses pages 5 à 9).
\item[5.] \emph{Propriétés formelles} (page 20) : (Hot~4) à (Hot~6).
\item[6.] \emph{Cofinalité} (page 21), \emph{recollement} (page 22),
\emph{la suite des lacets} (page 24), sur trois feuillets non paginés.
\end{enumerate}

\emph{Conventions.} Une \emph{théorie} $\mathcal H$ associe à toute petite
catégorie $A$ une catégorie $\mathcal H(A)$, et à tout foncteur $f : A \to B$
un foncteur $f^{*} : \mathcal H(B) \to \mathcal H(A)$, de façon
2-fonctorielle stricte et contravariante en $f$. On écrit $f^{*}$, $f_{!}$,
$f_{*}$ sans l'indice $\mathcal H$ que la page place en haut ou en bas.
$\Delta_0$ est la catégorie ponctuelle, $\mathrm{Hot}_{\mathcal H} =
\mathcal H(\Delta_0)$, et $\pi_A : A \to \Delta_0$ le foncteur unique. On
note $\int_A = (\pi_A)_{!}$ l'intégration et $\prod_A = (\pi_A)_{*}$ la
cointégration ; dans les exemples homotopiques, ce sont la limite inductive
et la limite projective \emph{homotopiques}. $A^{\wedge}$ est la catégorie
des préfaisceaux sur $A$ et $h_A : A \to A^{\wedge}$ le foncteur de Yoneda.
Les mots « cofibré » et « fibré » désignent les catégories cofibrées et
fibrées au sens de Grothendieck.\footnote{C'est ainsi qu'on lit les
« $f$ cofibrant » et « $g$ fibrant » de la page 20 ; la page ne définit pas
ces mots. Dans le vocabulaire actuel : op-fibration et fibration.}

\emph{Dates.} Deux feuillets sont datés de sa main : « déc. 83 » (page 1) et
« 24.11.83 » (page 4, en marge). Rien ne date les autres, et rien n'ordonne
les quatre suites entre elles.

\emph{Ce que le dossier annonce et n'établit pas} : les axiomes eux-mêmes
(« Dégager les bons axiomes ») ; l'existence des $f_{*}$ pour un
localisateur fondamental (« non encore ») ; la comparaison de deux couples
$(M, W)$ qui donnent la même théorie ; la condition de recollement de la
page 22 (« on veut ») ; le cas des catégories filtrantes (page 21, en marge).

\pagerange{1}{2}
\section*{La structure des dérivateurs (pages 1 et 2)}

Le feuillet 1, encadré sous ce titre, fixe une donnée $\mathcal H(I)$ pour
chaque petite catégorie $I$ et distingue trois niveaux de morphismes entre
deux objets $x, y$ de $\mathcal H(I)$ :
\begin{itemize}
  \item un ensemble $\mathrm{Hom}(x, y)$ ;
  \item un type d'homotopie $\mathrm{Homhot}(x, y) \in \mathrm{Hot}$, dont
  l'ensemble des composantes est $\mathrm{Hom}(x, y) = \pi_0\,
  \mathrm{Homhot}(x, y)$ ;
  \item un objet interne $\underline{\mathrm{Homhot}}(x, y) \in
  \mathrm{Hot}(I)$, dont on retrouve le précédent en passant au global sur
  $I$ :
  \[
    \mathrm{Homhot}(x, y) \simeq \mathrm{holim}_I\,
    \underline{\mathrm{Homhot}}(x, y).
  \]
\end{itemize}
\footnote{La page écrit $\int_I$. Or, dans la suite du dossier, $\int$ est
l'intégration, qui joue le rôle d'une limite \emph{inductive} (page 9). Un
espace de morphismes global s'obtient par une limite \emph{projective}
homotopique, la cointégration $\prod_I$ de la page 11 ; c'est elle qu'on
écrit. Le premier terme de la formule est surchargé sur la page, et la
nature de l'objet interne n'y est pas précisée.}

Pour $f : I \to J$, $X \in \mathcal H(I)$ et $Y \in \mathcal H(J)$,
l'adjonction $f_{!} \dashv f^{*}$ s'écrit aux trois niveaux :
\[
  \mathrm{Hom}(f_{!}X, Y) \simeq \mathrm{Hom}(X, f^{*}Y), \qquad
  \mathrm{Homhot}(f_{!}X, Y) \simeq \mathrm{Homhot}(X, f^{*}Y),
\]
\[
  \underline{\mathrm{Homhot}}(f_{!}X, Y) \simeq f_{*}\,
  \underline{\mathrm{Homhot}}(X, f^{*}Y),
\]
chacune donnant la précédente, la deuxième par passage aux $\pi_0$, la
troisième par passage au global sur $J$. De même, « duales », pour
l'adjonction $f^{*} \dashv f_{*}$ :
\[
  \mathrm{Hom}(f^{*}Y, X) \simeq \mathrm{Hom}(Y, f_{*}X), \qquad
  \mathrm{Homhot}(f^{*}Y, X) \simeq \mathrm{Homhot}(Y, f_{*}X),
\]
\[
  \underline{\mathrm{Homhot}}(Y, f_{*}X) \simeq f_{*}\,
  \underline{\mathrm{Homhot}}(f^{*}Y, X).
\]
Les deux formules internes sont encadrées.\footnote{Nous les donnons comme la
page les donne. Elles valent dans une théorie « enrichie » en types
d'homotopie, munie d'un $\mathrm{Hom}$ interne convenable ; la page ne dit pas
lequel. Pour le $\mathrm{Hom}$ interne cartésien des diagrammes d'espaces, la
première demanderait une formule de projection $f_{!}(X \times f^{*}Z)
\simeq f_{!}X \times Z$, qui n'a pas lieu pour tout $f$.}

Le feuillet 2 est un brouillon sans phrase : $\pi_0$ et $\pi_1$ d'un objet de
$\mathcal H(\Delta_1)$, des 2-flèches entre deux flèches $x
\rightrightarrows y$, une sous-catégorie $J \subset I$ et la restriction
$\mathcal H(I) \to \mathcal H(J)$, et, isolés, les mots « $M$ $W$ »,
« $\mathrm{Hot}$ », « $\mathrm{Cl}(M, W)$ ».

\pagerange{3}{4}
\section*{Un programme (pages 3 et 4)}

Deux feuillets énumèrent, sous des lettres a) à i), ce qu'une théorie
homotopique devrait expliquer :
\begin{enumerate}
  \item[a)] les \emph{modèles}, leur importance et leur indétermination ;
  \item[b)] les suites exactes longues associées à une flèche, liées aux
  cylindres et cocylindres ;
  \item[c)] les résolutions d'un modèle, à gauche ou à droite, et le
  formalisme des foncteurs dérivés quand deux théories sont en présence ;
  \item[d)] les notions homotopiques sur les modèles --- équivalences
  faibles $W$, cofibrations, homotopismes --- qu'une note marginale dit
  « intermédiaires » pour construire les catégories dérivées ;
  \item[e)] les homotopies entre applications, entre homotopies, et ainsi de
  suite, c'est-à-dire un espace $\mathrm{Hom}(X, Y)$ dans $\mathrm{Hot}$ ---
  « dans Quillen, il y a des $\mathrm{Hom}$ simpliciaux » ;
  \item[f)] des opérations, « grossières » sur les modèles ou « délicates »
  sur $\mathcal H$ ;
  \item[g)] des invariants $\pi_i$ ou $H_i$, à valeurs dans une catégorie
  souvent abélienne ;
  \item[h)] des invariants $K(\mathcal H)$ définis par « triangles
  exacts », et des $K_i$ supérieurs --- « cas additif seulement ?!? » ;
  \item[i)] des $\mathrm{Hom}$ internes et une structure multiplicative, avec
  \[
    \mathrm{Hom}(x \otimes y, z) \simeq \mathrm{Hom}(x,
    \underline{\mathrm{Hom}}(y, z)).
  \]
\end{enumerate}
\footnote{Les lettres des quatre derniers points sont surchargées ; f) et g)
sont restituées d'après la note marginale qui les nomme (« des aspects
particulièrement peu compris »), h) et i) d'après l'ordre. Une note datée du
24.11.83 demande, en marge du point i), les relations entre ces
$\mathrm{Hom}$ internes et les $\mathrm{Homhot}$ du point e).}

La conclusion : « sauf peut-être f) et g), il semble que tout puisse se
décrire en termes de $(M, W \subset \mathrm{Fl}(M))$ » --- une catégorie et
une classe de flèches à inverser. « Dégager les bons axiomes. »\footnote{Un
couple $(M, W)$ s'appelle aujourd'hui une catégorie relative, ou
homotopique. Que la localisation simpliciale d'un tel couple en retienne
toute la théorie homotopique est un théorème de Dwyer et Kan (1980) ; les
catégories relatives comme modèle des théories homotopiques sont étudiées
par Barwick et Kan (2012). Références citées de mémoire ; les pages ne
citent que Quillen.}

\pagerange{5}{11}
\section*{La variance et les premiers axiomes (pages 5 à 11)}

Le feuillet 5 ouvre une suite qu'il pagine de 1 à 9, sous le titre
« Théorie (de variance) homotopique $\mathcal H$ ».

\subsection*{La variance}

Une théorie $\mathcal H$ associe à toute petite catégorie $A$ une catégorie
$\mathcal H(A)$, non nécessairement petite, et à tout foncteur $f : A \to B$
un foncteur $f^{*} : \mathcal H(B) \to \mathcal H(A)$, de façon
2-contravariante \emph{stricte}. On pose $\mathrm{Hot}_{\mathcal H} =
\mathcal H(\Delta_0)$. Pour chaque objet $a$ de $A$, soit $i_a : \Delta_0 \to
A$ le foncteur de valeur $a$ ; on en tire le foncteur « diagramme sous-jacent »
\[
  \alpha_A : \mathcal H(A) \longrightarrow \underline{\mathrm{Hom}}(A,
  \mathrm{Hot}_{\mathcal H}), \qquad \mathcal X \longmapsto \bigl(a \mapsto
  i_a^{*}\mathcal X\bigr),
\]
naturel en $A$ : pour $f : A \to B$, $\alpha_A \circ f^{*} = f^{*} \circ
\alpha_B$.\footnote{Plus généralement, dit une N.B., tout $B$ donne
$\mathcal H(A) \to \underline{\mathrm{Hom}}(\underline{\mathrm{Hom}}(B, A),
\mathcal H(B))$, et $B = \Delta_0$ redonne $\alpha_A$. Dans le vocabulaire
actuel, une telle donnée est un \emph{prédérivateur} (Grothendieck,
\emph{Les Dérivateurs}, 1990--1991 ; Heller, \emph{Homotopy theories}, 1988 ;
cités de mémoire), et $\alpha_A$ le foncteur « diagramme ». La page n'impose
pas que la famille des $i_a^{*}$ soit conservative, qui est l'axiome
Der~2 des dérivateurs.}

\textbf{(Hot 1)} Si $A = \coprod_{i \in I} A_i$, alors $\mathcal H(A)
\to \prod_{i \in I} \mathcal H(A_i)$ est une équivalence.\footnote{C'est
l'axiome Der~1.}

\subsection*{Une première intégration, abandonnée (page 7)}

Un premier essai définit l'intégration en passant par les préfaisceaux : on
se donne, pour toute petite $A$, un foncteur $L_A : \mathcal H(A) \to
\mathcal H(A^{\wedge})$ tel que $(h_A)^{*} L_A = \mathrm{id}$, et l'on
demande que les $f_{!}$ soient induits par un carré commutatif avec
$L_A$, $L_B$ et le foncteur induit $\mathcal H(A^{\wedge}) \to
\mathcal H(B^{\wedge})$. Tout ce passage est barré de longs traits
obliques.\footnote{Il suppose $\mathcal H$ défini sur $A^{\wedge}$, qui
n'est pas petite ; les « $U$-catégories » de la page 5 laissent penser à un
changement d'univers. La page 21 reprend le même carré, avec $(h_A)_{!}$ à
la place de $L_A$, pour un usage précis.}

\subsection*{Les adjoints (pages 9 et 11)}

\textbf{(Hot 2)} Pour tout $f : A \to B$, le foncteur $f^{*} : \mathcal H(B)
\to \mathcal H(A)$ a un adjoint à gauche $f_{!} : \mathcal H(A) \to
\mathcal H(B)$.\footnote{La page écrit $f^{\mathcal H}_{!} : \mathcal H(B) \to
\mathcal H(A)$, la source et le but de $f^{*}$ ; c'est une inadvertance.
Une remarque, « A. », se demande s'il ne faudra pas se borner aux $f$
« cofibrantes ». Les axiomes actuels (Der~3) ne se restreignent pas.}

Pour $B = \Delta_0$, on obtient l'\emph{intégration} $\int_A =
(\pi_A)_{!} : \mathcal H(A) \to \mathrm{Hot}_{\mathcal H}$, qui joue le rôle de
$\varinjlim_A$. Par transitivité, $\int_A \mathcal X \simeq \int_B
f_{!}\mathcal X$. L'existence des $f_{!}$ implique celle des sommes dans les
$\mathcal H(A)$, avec (Hot~1) ; et les $f^{*}$ y commutent dès qu'ils ont
aussi un adjoint à droite.\footnote{La page dit « et les $f^{*}$ bien sûr y
commutent ». C'est vrai sous (Hot~3), qui fait de $f^{*}$ un adjoint à
gauche ; sous (Hot~2) seul, il faudrait une propriété de changement de base,
celle de (Hot~4).}

\textbf{(Hot 3)} Pour tout $f$, le foncteur $f^{*}$ a un adjoint à droite
$f_{*}$. D'où la \emph{cointégration} $\prod_A = (\pi_A)_{*}$, avec $\prod_A
\mathcal X \simeq \prod_B f_{*}\mathcal X$, et l'existence des produits dans
les $\mathcal H(A)$.\footnote{(Hot~2) et (Hot~3) forment l'axiome Der~3 :
existence des extensions de Kan homotopiques à gauche et à droite.}

\emph{Dualité.} En posant $\mathcal H^{\circ}(A) = \mathcal H(A)^{\circ}$,
on obtient une théorie $\mathcal H^{\circ}$ où $f^{*}_{\mathcal H^{\circ}} =
(f^{*})^{\circ}$, $f^{\mathcal H^{\circ}}_{!} = (f_{*})^{\circ}$ et
$f^{\mathcal H^{\circ}}_{*} = (f_{!})^{\circ}$ : intégration et cointégration
s'échangent.

\emph{Cas additif.} Si les $\mathcal H(A)$ sont additives (ou $k$-linéaires)
et les $f^{*}$ additifs, les $f_{!}$ et $f_{*}$ le sont aussi, comme adjoints
de foncteurs additifs.

\pagerange{13}{20}
\section*{Quatre exemples (pages 13 à 20)}

\emph{Exemple 1.} Soit $M$ une catégorie complète et cocomplète, et
$\mathcal H(A) = \underline{\mathrm{Hom}}(A, M)$. Les $f^{*}$ sont les
restrictions, leurs adjoints les extensions de Kan à gauche et à droite ;
$\mathrm{Hot}_{\mathcal H} = M$, $\int_A = \varinjlim_A$ et $\prod_A =
\varprojlim_A$. C'est le dérivateur représenté par $M$.

\emph{Exemple 2.} Soit $(M, W)$ une catégorie munie d'une classe de flèches
$W$. On note $W_A$ la classe des flèches de $\underline{\mathrm{Hom}}(A, M)$
qui sont dans $W$ objet par objet, et l'on pose
\[
  \mathcal H(A) = W_A^{-1}\,\underline{\mathrm{Hom}}(A, M).
\]
Les $f^{*}$ passent aux localisées, puisqu'ils préservent les $W_A$.
L'existence des adjoints est une condition sur $W$.\footnote{Elle est
satisfaite, par exemple, quand $(M, W)$ provient d'une catégorie de modèles
de Quillen : c'est le théorème selon lequel toute catégorie de modèles
définit un dérivateur (Cisinski, 2003, cité de mémoire). La note marginale
affirme que (Hot~1) est satisfait. Ce n'est pas automatique : une
localisation commute aux produits finis, pas en général aux produits
infinis ; pour une catégorie de modèles, cela vaut.}

\emph{Une remarque sur les morphismes} (pages 15 et 17). Un foncteur $f : M \to
M'$ tel que $f(W) \subset W'$ induit des foncteurs $\mathcal H_W(A) \to
\mathcal H_{W'}(A)$ compatibles aux $f^{*}$, mais rien ne dit qu'ils
commutent aux $f_{!}$ et $f_{*}$ : quand $W$ et $W'$ sont réduits aux
isomorphismes, cela voudrait dire que $f$ commute aux limites inductives,
resp.\ projectives. En revanche, si ce sont des équivalences pour tout $A$,
la commutation suit. En voici un cas. Supposons que $f$ ait un adjoint à
droite $g$, que $f^{-1}(W') = W$, que la counité $fg(X') \to X'$ soit dans
$W'$ pour tout $X'$, et que $W'$ satisfasse au « deux sur trois ». Alors $g$
envoie $W'$ dans $W$, l'unité et la counité sont dans $W$ et $W'$, et $f$ et
$g$ induisent des équivalences inverses $\mathcal H_W(A) \simeq
\mathcal H_{W'}(A)$ pour tout $A$, les mêmes conditions passant à
$\underline{\mathrm{Hom}}(A, M) \rightleftarrows \underline{\mathrm{Hom}}(A,
M')$ objet par objet.\footnote{La page n'écrit pas la propriété « deux sur
trois », sans laquelle l'argument ne passe pas. Avec elle : si $w' : X' \to
Y'$ est dans $W'$, alors $fg(w')$ l'est, car la counité est naturelle, donc
$g(w') \in f^{-1}(W') = W$ ; et $f$ appliqué à l'unité, composé avec la
counité en $f(X)$, donne l'identité, donc l'unité est dans $W$. La
conclusion de la page est barrée après « Donc » ; nous l'écrivons.}

Les théories qui se décrivent par un couple $(M, W)$ sont dites
\emph{spéciales}, et « cela implique des conditions draconiennes sur $(M,
W)$ ». On peut se proposer d'avoir une vue d'ensemble des couples qui
donnent une théorie donnée : deux tels couples $(M_1, W_1)$, $(M_2, W_2)$
s'envoient-ils toujours dans un troisième, de façon à induire l'isomorphisme
donné entre les théories ?\footnote{La question reste ouverte sur la page.
Aujourd'hui, pour les catégories de modèles combinatoires, deux catégories
qui définissent des dérivateurs équivalents sont reliées par une chaîne
d'équivalences de Quillen (Renaudin, 2009, cité de mémoire) ; hors de ce
cadre, la réponse dépend des hypothèses.}

\emph{Exemple 3.} $M = (\mathrm{Cat})$ et $W$ un \emph{localisateur
fondamental}. La page dit : « J'ai prouvé l'existence des $f_{!}$, mais non
encore celle des $f_{*}$ », et ajoute qu'il faudra peut-être des hypothèses
de finitude sur $A$, $B$.\footnote{La page ne donne pas la démonstration.
Que tout localisateur fondamental définisse un dérivateur, avec les deux
adjoints et sans hypothèse de finitude, est établi par Cisinski
(\emph{Les préfaisceaux comme modèles des types d'homotopie}, Astérisque
308, 2006), cité de mémoire. Le terme « localisateur fondamental » est aussi
celui de \emph{À la poursuite des champs}, dont le dossier relève dans
l'inventaire.}

\emph{Exemple 4.} Soit $\mathcal A$ une catégorie abélienne, $M$ ses
complexes et $W$ les quasi-isomorphismes : $\mathcal H(A)$ est la catégorie
dérivée de $\underline{\mathrm{Hom}}(A, \mathcal A)$. Si $\mathcal A$ est la
catégorie des faisceaux de $k$-modules sur un topos $\mathcal T$, alors
$\underline{\mathrm{Hom}}(A, \mathcal A)$ est la catégorie des faisceaux de
$k$-modules sur le topos produit $A^{\circ\wedge} \times \mathcal T$, et
$f^{*}$ est l'image inverse par le morphisme de topos $f^{\circ} \times
\mathrm{id}_{\mathcal T}$. Le foncteur induit sur les catégories dérivées a
un adjoint à droite $\mathrm{R}f_{*}$, et aussi un adjoint à gauche
$\mathrm{L}f_{!}$.\footnote{La page dit d'abord que, « généralement », il
n'y a pas d'adjoint à gauche, sauf pour le topos ponctuel. Puis elle obtient
$\mathrm{L}f_{!}$ ainsi : au niveau des faisceaux d'ensembles, $(f^{\circ})_{!}$
existe et serait « exact à gauche », donc transformerait quasi-isomorphismes
en quasi-isomorphismes. L'argument ne tient pas tel qu'il est écrit : sur
les faisceaux abéliens, l'extension de Kan à gauche est exacte à droite, non
exacte, et ne préserve pas les quasi-isomorphismes. La conclusion est juste
par une autre voie : les complexes d'une catégorie abélienne de Grothendieck
définissent un dérivateur, avec les deux adjoints (Cisinski, 2003, cité de
mémoire), et $\mathrm{L}f_{!}$ se calcule par des résolutions convenables.}

\pagerange{20}{20}
\section*{Propriétés formelles (page 20)}

\textbf{(Hot 4) Changement de base.} Soit un carré cartésien de petites
catégories
\[
  A = B \times_{B'} A' \xrightarrow{\ g'\ } A', \qquad f : A \to B, \qquad
  f' : A' \to B', \qquad g : B \to B'.
\]
Les morphismes canoniques
\[
  f_{!}\, g'^{*} \longrightarrow g^{*} f'_{!} \qquad \text{et} \qquad
  g^{*} f'_{*} \longrightarrow f_{*}\, g'^{*}
\]
sont des isomorphismes : le premier si $f'$ est cofibré, ou si $g$ est
fibré ; le second si $f'$ est fibré, ou si $g$ est cofibré.\footnote{La page
met la condition sur $f$, qui dans un carré cartésien l'hérite de $f'$ ; c'est
$f'$ qui compte. Elle ajoute les alternatives entre crochets, « [ou $g$ fibrant
??] » et « [ou $g$ cofibrant ??] », avec ses points d'interrogation. Elles
sont justes : un carré cartésien est homotopiquement exact dès que sa flèche
verticale est propre, par exemple cofibrée, ou sa flèche horizontale lisse,
par exemple fibrée, et dualement (Grothendieck, \emph{Les Dérivateurs} ;
Maltsiniotis, 2005 ; cités de mémoire). Les axiomes actuels (Der~4)
demandent le changement de base pour les carrés « comma », dont celui-ci
se déduit. La page note enfin que les deux formules sont « transposées »
l'une de l'autre.}

\textbf{(Hot 5)} Si $f : A \to B$ est pleinement fidèle, les foncteurs
$f_{!}$ et $f_{*}$ sont pleinement fidèles : $f^{*}f_{!} \simeq \mathrm{id}$ et
$f^{*}f_{*} \simeq \mathrm{id}$. Autrement dit, $f^{*}$ est un foncteur de
localisation.

\textbf{(Hot 6)} Si $f : A \to B$ est une localisation \emph{au sens
homotopique}, alors $f^{*}$ est pleinement fidèle : $f_{!}f^{*} \simeq
\mathrm{id} \simeq f_{*}f^{*}$.\footnote{La page dit « si $f$ localisation »,
sans préciser. Pour une localisation au sens ordinaire des catégories,
l'énoncé est faux. Soit $A$ une catégorie dont le nerf a le type
d'homotopie de la sphère $S^{2}$ --- l'ensemble ordonné des faces propres d'un
tétraèdre, par exemple --- et $f : A \to \Delta_0$. Comme $A$ est connexe et
simplement connexe, localiser $A$ en toutes ses flèches donne une catégorie
équivalente à $\Delta_0$ : $f$ est une localisation. Mais, dans la théorie
des types d'homotopie, $\mathrm{Hom}(f^{*}X, f^{*}Y)$ est l'ensemble des
classes d'applications $X \times S^{2} \to Y$, qui diffère de
$\mathrm{Hom}(X, Y)$ pour $X$ ponctuel et $Y = K(\mathbb Z, 2)$. L'énoncé
vaut quand le nerf de $f$ est une localisation homotopique, au sens de
Dwyer et Kan. Le contre-exemple est le nôtre ; la référence est citée de
mémoire. La dernière formule est écrite contre le bord du feuillet.}

\pagerange{21}{21}
\section*{Cofinalité (page 21)}

On cherche des conditions sur un foncteur $i : A_0 \to A$ pour que
\[
  \int_{A_0} i^{*}\xi \xrightarrow{\ \sim\ } \int_A \xi \qquad \text{pour tout
  } \xi \in \mathcal H(A),
\]
ce qu'il appelle « cofinalité ». Il en donne une : $i$ a un adjoint à gauche
$j : A \to A_0$. Alors
\[
  i^{*} \simeq j_{!},
\]
d'où, puisque $p_0 j = p$ ($p$ et $p_0$ désignant les foncteurs vers
$\Delta_0$),
\[
  \int_{A_0} i^{*} = p_{0!}\, i^{*} \simeq p_{0!}\, j_{!} \simeq p_{!} =
  \int_A.
\]
\footnote{La page écrit la formule $i^{*} = j_{!}$ entourée, précédée d'un
point d'interrogation (« je présume »), puis les deux formules suivantes avec
$i$ là où l'argument demande $j$ : $p_0$ et $i$ ne se composent pas en
$p_0 i$. Nous écrivons $j$. L'énoncé est le fait classique qu'un adjoint à
droite est homotopiquement cofinal.}

\emph{Exemple.} Si $A$ a un objet final $e$ et que $A_0$ est un point envoyé
sur $e$, le foncteur $A \to A_0$ est adjoint à gauche de $i$, et $\int_A \xi
\simeq \xi(e)$.

\emph{Justification de $i^{*} \simeq j_{!}$.} Pour tout foncteur $j : A \to
B$, le carré
\[
  (h_B)_{!}\, j_{!} \simeq (j^{*}_{\wedge})^{*}\, (h_A)_{!}
\]
commute, où $j^{*}_{\wedge} : B^{\wedge} \to A^{\wedge}$ est la restriction des
préfaisceaux.\footnote{Les deux membres sont des extensions de Kan à gauche
calculées sur la même catégorie comma, celle des éléments de $j^{*}G$ pour
$G \in B^{\wedge}$ ; l'égalité suit du changement de base (Hot~4). Ce calcul
est le nôtre ; la page renvoie à « un diagramme commutatif ».} Si $j$ a un
adjoint à droite $i$, on a $j^{*}_{\wedge} = i^{\wedge}_{!}$, qui prolonge
$i$ aux préfaisceaux. En restreignant par $h_B$, pleinement fidèle, et en
utilisant (Hot~5), on obtient $j_{!} \simeq i^{*}$.

En marge : « Cas des catégories filtrantes ?? ».\footnote{Question sans
suite. Pour les dérivateurs, une limite inductive homotopique filtrante n'est
pas en général calculée objet par objet ; c'est une propriété
supplémentaire de certains dérivateurs.}

\pagerange{22}{22}
\section*{Un recollement (page 22)}

Soient $A'$, $A''$ deux sous-catégories pleines de $A$, toutes deux des
cribles ou toutes deux des cocribles, de réunion $A$ et d'intersection $A_0$.
Alors $A = A' \amalg_{A_0} A''$ dans $(\mathrm{Cat})$, puisque toute flèche de
$A$ est dans $A'$ ou dans $A''$. On voudrait que le foncteur
\[
  \mathcal H(A) \longrightarrow \mathcal H(A') \times_{\mathcal H(A_0)}
  \mathcal H(A'')
\]
soit \emph{plein} --- pour $\xi, \eta \in \mathcal H(A)$, l'application
$\mathrm{Hom}(\xi, \eta) \to \mathrm{Hom}(\xi', \eta') \times_{\mathrm{Hom}(\xi_0,
\eta_0)} \mathrm{Hom}(\xi'', \eta'')$ surjective --- et il « a tendance d'être
bijectif sur les objets ». On s'en servira pour la catégorie en forme de coin
$a' \leftarrow c \to a''$, réunion de deux flèches issues de $c$, pour son
dual, et pour leurs produits par une petite catégorie $B$.\footnote{C'est une
forme de l'axiome dit « fort » des dérivateurs, qui demande que
$\mathcal H(B \times \Delta_1) \to \mathcal H(B)^{\Delta_1}$ soit plein et
essentiellement surjectif (Der~5 ; Heller ; cités de mémoire). La
condition de la page, pour les réunions de deux cribles, en est voisine ;
nous ne l'identifions pas.}

\pagerange{24}{24}
\section*{La suite des lacets (page 24)}

Le feuillet 24 ne porte que des diagrammes, dans une théorie pointée, $e$
désignant l'objet nul. Pour une flèche $f : X \to Y$, on complète par des
carrés homotopiquement cartésiens :
\[
  Z = X \times^{h}_{Y} e, \qquad \Omega Y = e \times^{h}_{Y} e, \qquad
  \Omega Y \xrightarrow{\ h\ } Z,
\]
et l'on recommence. Les carrés s'empilent en escalier jusqu'à
$\Omega^{4} Y$, et la suite
\[
  Y \xleftarrow{\ f\ } X \xleftarrow{\ g\ } Z \xleftarrow{\ h\ } \Omega Y
  \xleftarrow{\ \Omega f\ } \Omega X \xleftarrow{\ \Omega g\ } \Omega Z
  \xleftarrow{\ \Omega h\ } \Omega^{2} Y \longleftarrow \cdots
\]
est redessinée en zigzag, chaque composé de deux flèches consécutives passant
par $e$. C'est la suite des fibres, ou suite de Puppe.\footnote{La page ne
dit pas que les carrés sont homotopiquement cartésiens ; c'est ce qui donne à
l'escalier son sens, et nous le précisons. Une telle construction de la suite
des fibres par carrés cartésiens itérés est standard dans les dérivateurs
pointés (Groth, 2013, cité de mémoire). Le feuillet porte encore un croquis
en perspective de l'escalier, un carré en pointillés et un petit croquis à
deux lettres, que nous ne reconstituons pas. Les signes de la flèche $\Omega
h$ et le détail des étiquettes sont ceux de la page ; les carrés composés
disent que $\Omega$ de la suite est, au signe près, la suite décalée.}

\end{document}
